\begin{figure}[H]\centering\resizebox{0.98\linewidth}{!}{%
\begin{tikzpicture}[
 b/.style={draw,rounded corners=1.2pt,align=center,minimum width=2.8cm,minimum height=0.9cm,fill=blue!6},
 p/.style={draw,rounded corners=1.2pt,align=center,minimum width=3.4cm,minimum height=0.9cm,fill=green!7},
 a/.style={-Stealth,thick},font=\small]
\node[b] (source) at (0,0) {普通 C 源码\\保持测试程序源文件不变};
\node[b] (gimple) at (4.1,0) {GIMPLE\\语言无关中间表示};
\node[p] (shape) at (8.7,1.35) {跨基本块结构识别\\循环、链式更新与归约};
\node[p] (rtl) at (8.7,-1.35) {RTL 局部模式识别\\算术与位操作序列};
\node[b,fill=cyan!7] (ifn) at (12.9,1.35) {internal function\\保留明确硬件语义};
\node[b,fill=orange!8] (insn) at (17.0,0) {机器描述 \code{define\_insn}\\约束、调度属性与编码};
\node[b,fill=gray!10] (code) at (21.2,0) {目标机器码\\自动生成自定义指令};

\draw[a] (source)--(gimple);
\draw[a] ([yshift=0.18cm]gimple.east) -- (shape.west);
\draw[a] ([yshift=-0.18cm]gimple.east) -- (rtl.west);
\draw[a] (shape)--(ifn);
\draw[a] (ifn.east) -- ([yshift=0.2cm]insn.west);
\draw[a] (rtl.east) -- node[pos=0.53,below]{\code{peephole2}} ([yshift=-0.2cm]insn.west);
\draw[a] (insn)--(code);
\end{tikzpicture}}
\caption{从普通 C 程序语义到厂商自定义指令的 GCC 自动选择流程}
\label{fig:compiler-selection}
\end{figure}
